\chapter{Shifted QR}
\label{ch:shifted-qr}

Supplementary reading:
\begin{itemize}
  \item \cite{Demmel} Section 4.4.8 (Theorem 4.9, the implicit Q theorem).
  \item \cite{Demmel} Sections 4.4.4 and 4.4.5 (QR iteration with a shift, making it practical).
  \item \cite{Demmel} Section 5.3.1 (Wilkinson's shift).
  \item \cite{TB} Lecture 33 (Arnoldi iteration).
\end{itemize}

\section{Recap: QR Iteration on Hessenberg Matrices}

The lecture opens with a recap of \Cref{ch:hessenberg-qr}. We first reduce $A$ once to upper Hessenberg form $H = Q^\T AQ$ (\Cref{sec:hessenberg}). Each step of the QR iteration on $T_0 = H$ then consists of a QR factorization by Givens rotations followed by the reverse product (\Cref{sec:qr-hessenberg}),
\begin{equation}
  R_k = G_{m-1}^\T\cdots G_1^\T\,T_{k-1}, \qquad T_k = R_kG_1\cdots G_{m-1},
\end{equation}
and costs $O(m^2)$ because $T_k$ is again upper Hessenberg.

The board compares two orders of applying the same rotations. In the \emph{two-sweep} order all left multiplications are done before any right multiplication,
\begin{equation}
  G_1^\T T_{k-1}, \quad G_2^\T G_1^\T T_{k-1}, \quad \dots, \quad R_k, \quad R_kG_1, \quad R_kG_1G_2, \quad \dots, \quad T_k,
\end{equation}
while \emph{bulge chasing} (\Cref{sec:bulge-chasing}) applies every rotation from both sides as soon as it is known,
\begin{equation}
  G_1^\T T_{k-1}G_1, \quad G_2^\T G_1^\T T_{k-1}G_1G_2, \quad \dots, \quad T_k.
\end{equation}
By associativity both orders compute the same product \emph{if they use the same rotations}. But in bulge chasing, $G_2, \dots, G_{m-1}$ are chosen to remove the bulge that $G_1^\T T_{k-1}G_1$ creates at $(3,1)$, not to remove the subdiagonal entries of $G_1^\T T_{k-1}$. Why does this give the same $T_k$? The answer is the implicit Q theorem, stated at the end of \Cref{ch:hessenberg-qr}; this lecture proves it, then turns to shifts.

\section{Proof of the Implicit Q Theorem}
\label{sec:implicit-q-proof}

Recall the statement (\Cref{thm:implicit-q}): if $Q$ is orthogonal and $Q^\T AQ = H$ is unreduced upper Hessenberg, then columns $2, \dots, m$ of $Q$ are uniquely determined by the first column of $Q$, up to signs.

\begin{proof}
  We follow~\cite[Theorem 4.9]{Demmel}. Let both $Q^\T AQ = H$ and $V^\T AV = G$ be as in the theorem, with equal first columns $\vec{q}_1 = \vec{v}_1$. We want to show that
  \begin{equation}
    W \defeq V^\T Q = \diag(\pm 1, \dots, \pm 1),
  \end{equation}
  since then $Q = VW$, i.e., $\vec{q}_i = \pm\vec{v}_i$ for every $i$.

  $W$ is orthogonal as a product of orthogonal matrices, and its first column is $W\vec{e}_1 = V^\T\vec{q}_1 = V^\T\vec{v}_1 = \vec{e}_1$. Moreover, using $GV^\T = V^\T AVV^\T = V^\T A$ and $AQ = QH$,
  \begin{equation}
    GW = GV^\T Q = V^\T AQ = V^\T QH = WH.
  \end{equation}
  Write $\vec{w}_i$ for the $i$th column of $W$. Taking column $i$ of both sides, and using that column $i$ of the upper Hessenberg matrix $H$ is zero below row $i + 1$,
  \begin{equation}
    G\vec{w}_i = (GW)_{:,i} = (WH)_{:,i} = \sum_{j=1}^{i+1}\vec{w}_jH_{ji}.
  \end{equation}
  The term $j = i + 1$ is the only one that contains $\vec{w}_{i+1}$. Moving the others to the right-hand side gives
  \begin{equation}
    \label{eq:implicit-q-recursion}
    H_{i+1,i}\,\vec{w}_{i+1} = G\vec{w}_i - \sum_{j=1}^{i}H_{ji}\vec{w}_j.
  \end{equation}
  We show by induction that each $\vec{w}_i$ is nonzero in entries $1, \dots, i$ only:
  \begin{equation}
    \vec{w}_i = \begin{bmatrix} \times \\ \vdots \\ \times \\ 0 \\ \vdots \\ 0 \end{bmatrix} \begin{matrix} \\ \\ \leftarrow i \\ \\ \\ \\ \end{matrix}
  \end{equation}
  This holds for $\vec{w}_1 = \vec{e}_1$. Suppose it holds for $\vec{w}_1, \dots, \vec{w}_i$. Then $G\vec{w}_i$ is a combination of the first $i$ columns of $G$, which are zero below row $i + 1$ because $G$ is upper Hessenberg. The sum in \eqref{eq:implicit-q-recursion} is zero below row $i$ by the induction hypothesis. So the right-hand side of \eqref{eq:implicit-q-recursion} is zero below row $i + 1$. Since $H$ is unreduced, $H_{i+1,i} \ne 0$, and dividing by it shows that $\vec{w}_{i+1}$ is zero below row $i + 1$.

  Hence $W$ is upper triangular. Since $W$ is also orthogonal, $W^{-1} = W^\T$ is both upper triangular (as the inverse of an upper triangular matrix) and lower triangular (as the transpose of one), so $W$ is diagonal. An orthogonal diagonal matrix has $w_{ii}^2 = 1$, so $W = \diag(\pm 1, \dots, \pm 1)$.
\end{proof}

The proof has the same structure as the proof that the QR factorization is unique (\Cref{thm:qr-unique}). We take two solutions and look at the matrix $W$ that converts one into the other. Showing that $W$ is triangular and orthogonal forces it to be diagonal with entries of modulus one, and these are exactly the sign choices that cannot be fixed.

The recursion \eqref{eq:implicit-q-recursion} is also a constructive version of the theorem. Given $\vec{q}_1$, the relation $AQ = QH$ determines $h_{1i}, \dots, h_{ii}$ by orthogonality and then $\vec{q}_{i+1}$ by normalization, one column at a time. This is the Arnoldi iteration~\cite[Lecture 33]{TB}, and the step that divides by $h_{i+1,i}$ is where unreducedness is needed.

For $\K = \C$ the same proof works with $Q$ and $V$ unitary and $\T$ replaced by $\ast$. Then $W$ is a unitary diagonal matrix, so each column of $Q$ is determined up to a factor of modulus one.

\section{Application to Bulge Chasing}
\label{sec:implicit-q-bulge}

Both the two-sweep algorithm and bulge chasing compute $T_k = U_k^\T T_{k-1}U_k$ with $U_k$ a product of rotations $G_1G_2\cdots G_{m-1}$, where $G_j$ acts on rows $j$ and $j+1$. Only $G_1$ touches the first row, so $U_k\vec{e}_1 = G_1\vec{e}_1$ in both cases. Moreover, $G_1$ is the same in both, since it is determined by the first column of $T_{k-1}$ alone. Both results are upper Hessenberg. So if $T_k$ is unreduced (otherwise we deflate), \Cref{thm:implicit-q} shows that the two $U_k$ agree up to the signs of their columns. The two $T_k$ then agree up to a similarity $D^\T T_kD$ with $D = \diag(\pm 1, \dots, \pm 1)$, which only flips the signs of some rows and columns. As in \Cref{sec:basic-qr}, such a similarity leaves the diagonal and the moduli of all entries unchanged.

\section{The Shifted QR Iteration}
\label{sec:shifted-qr}

The QR iteration of \Cref{alg:basic-qr} can converge slowly, at the rate $\abs{\lambda_{p+1}/\lambda_p}$ (\Cref{sec:orth-iter-conv}). As for inverse iteration (\Cref{alg:inverse}), the remedy is a shift. We compare the two iterations side by side:
\begin{equation}
  \text{normal QR:}\quad
  \begin{aligned}
    T_{k-1} &= U_kR_k, \\
    T_k &= R_kU_k,
  \end{aligned}
  \qquad\qquad
  \text{shifted QR:}\quad
  \begin{aligned}
    T_{k-1} - \mu_kI &= U_kR_k, \\
    T_k &= R_kU_k + \mu_kI.
  \end{aligned}
\end{equation}
The shift $\mu_k$ is subtracted before the factorization and added back afterwards. We check the two properties that made the unshifted iteration work.

\paragraph{The upper Hessenberg form is preserved.} Subtracting and adding $\mu_kI$ changes only the diagonal. So $T_{k-1} - \mu_kI$ is upper Hessenberg, its QR factorization takes $m - 1$ Givens rotations, and $R_kU_k$ is upper Hessenberg by \Cref{sec:qr-hessenberg}. Each step still costs $O(m^2)$.

\paragraph{The eigenvalues are preserved.} From the first equation, $R_k = U_k^\ast(T_{k-1} - \mu_kI)$, so
\begin{equation}
  T_k = R_kU_k + \mu_kI = U_k^\ast(T_{k-1} - \mu_kI)U_k + \mu_kI = U_k^\ast T_{k-1}U_k,
\end{equation}
since $U_k^\ast(\mu_kI)U_k = \mu_kI$. So $T_k$ is again a unitary similarity transformation of $T_{k-1}$. Without adding back $\mu_kI$, the eigenvalues would be shifted by $-\mu_k$ in every step.

\paragraph{Tracking the eigenvectors.} Since $T_k = U_k^\ast T_{k-1}U_k$ has exactly the same form as in the unshifted iteration, we can track $Q_k$ as before (\Cref{sec:basic-qr}): accumulate $Q \gets QU_k$ after each step, so that $T_k = (U_1\cdots U_k)^\ast H(U_1\cdots U_k)$. Combined with the Hessenberg reduction, $A \approx (QU_1\cdots U_k)\,T_k\,(QU_1\cdots U_k)^\ast$ once $T_k$ is nearly triangular.

\section{Choosing the Shift}
\label{sec:choosing-shift}

How does the shift affect convergence, and which $\mu_k$ should we choose? We first consider the ideal case.

\paragraph{Exact shift.} Assume $\mu_k$ is an eigenvalue of $T_{k-1}$. Then $T_{k-1} - \mu_kI$ is singular. Assume moreover that all off-diagonal entries $t_{j+1,j}$ are nonzero, i.e., $T_{k-1}$ is unreduced (otherwise we split the problem, see below). Then the first $m - 1$ columns of $T_{k-1} - \mu_kI$ are linearly independent: rows $2, \dots, m$ of these columns form an upper triangular $(m-1) \times (m-1)$ matrix whose diagonal entries are the nonzero subdiagonal entries of $T_{k-1}$. In the QR factorization $T_{k-1} - \mu_kI = U_kR_k$, column $j$ is a combination of $\vec{u}_1, \dots, \vec{u}_j$, with the coefficient $r_{jj}$ on $\vec{u}_j$. If $r_{jj} = 0$ for some $j \le m - 1$, column $j$ would lie in the span of the previous columns. Hence $r_{11}, \dots, r_{m-1,m-1} \ne 0$, and since $R_k$ is singular, \emph{the last entry must be zero}, $r_{mm} = 0$. As $R_k$ is upper triangular, its whole last row vanishes:
\begin{equation}
  \underbrace{T_{k-1} - \mu_kI}_{\text{unreduced Hessenberg}} = U_k\,\underbrace{\begin{bmatrix} \times & \times & \times & \times \\ 0 & \times & \times & \times \\ 0 & 0 & \times & \times \\ 0 & 0 & 0 & 0 \end{bmatrix}}_{R_k}.
\end{equation}

Now what happens to $T_k$? The last row of $R_kU_k$ is the last row of $R_k$ times $U_k$, which is zero. Adding $\mu_kI$ gives
\begin{equation}
  T_k = R_kU_k + \mu_kI = \begin{bmatrix} \hat{T} & \vec{t} \\ 0 \;\cdots\; 0 & \mu_k \end{bmatrix},
\end{equation}
with $\hat{T} \in \K^{(m-1) \times (m-1)}$ upper Hessenberg. So $t_{m,m-1} = 0$ and $t_{mm} = \mu_k$: with an exact shift, \emph{one step} of the QR iteration splits off the eigenvalue $\mu_k$. In the words of the board, we don't have to worry about the last entry any more.

We then continue on $\hat{T}$ alone. If a unitary $Q$ acts on $\hat{T}$, then on the full matrix it acts as $\diag(Q, 1)$,
\begin{equation}
  \begin{bmatrix} Q^\ast & 0 \\ 0 & 1 \end{bmatrix}
  \begin{bmatrix} \hat{T} & \vec{t} \\ 0 & \mu_k \end{bmatrix}
  \begin{bmatrix} Q & 0 \\ 0 & 1 \end{bmatrix}
  = \begin{bmatrix} Q^\ast\hat{T}Q & Q^\ast\vec{t} \\ 0 & \mu_k \end{bmatrix},
\end{equation}
so the eigenvalue already found is left untouched, and the problem size has decreased by one. The last column $Q^\ast\vec{t}$ only needs to be updated if we want the full Schur form, not just the eigenvalues.

\paragraph{Splitting.} In practice $\mu_k$ is only close to an eigenvalue, and the subdiagonal entries become small rather than zero. Whenever a subdiagonal entry approaches $0$, we set it to zero; $T_k$ is then (approximately) block upper triangular,
\begin{equation}
  T_k \approx \begin{bmatrix} (T_k)_{11} & (T_k)_{12} \\ 0 & (T_k)_{22} \end{bmatrix},
\end{equation}
and its eigenvalues are those of the two diagonal blocks, which are again upper Hessenberg. We recurse on $(T_k)_{11}$ and $(T_k)_{22}$ separately.

\begin{algorithm}[H]
  \caption{Practical QR iteration.}
  \label{alg:practical-qr}
  \begin{algorithmic}[1]
    \Require $A \in \K^{m \times m}$
    \State compute $H = Q^\ast AQ$ \Comment{\Cref{alg:hessenberg}, once}
    \State $T_0 \gets H$, $k \gets 0$
    \While{not converged}
      \State $k \gets k + 1$, choose $\mu_k$
      \State $T_{k-1} - \mu_kI = U_kR_k$ \Comment{Givens QR, $O(m^2)$}
      \State $T_k \gets R_kU_k + \mu_kI$
      \If{some subdiagonal entry of $T_k$ is $\approx 0$}
        \State split $T_k$ into $(T_k)_{11}$ and $(T_k)_{22}$ and recurse on both
      \EndIf
    \EndWhile
  \end{algorithmic}
\end{algorithm}

\paragraph{Which eigenvalue to target?} The goal is to choose $\mu_k \approx$ an eigenvalue and to split up the problem wherever a subdiagonal entry is $\approx 0$. But which eigenvalue? The answer on the board is $\mu_k \approx \lambda_m$, for two reasons.
\begin{enumerate}
  \item The shift helps the eigenvalue $\lambda_i$ to separate if $\abs{\lambda_i - \mu} \ll \abs{\lambda_j - \mu}$ for all $j \ne i$. Without a shift ($\mu = 0$) this is easiest if $\lambda_i$ itself is small, i.e., for the eigenvalue of smallest modulus.
  \item The QR iteration orders the eigenvalues by size along the diagonal of $T_k$ (\Cref{thm:orth-iter-schur}), with the smallest one, $\lambda_m$, at the bottom right. A shift close to a different eigenvalue would have to move that eigenvalue to the bottom, and \emph{reordering the eigenvalues slows down convergence}.
\end{enumerate}
So the shift should approximate the eigenvalue that is already converging in the bottom right corner, where we then look for $t_{m,m-1} \approx 0$.

\section{Relationship with (Inverse) Power Iteration}
\label{sec:shifted-orth-iter}

We have seen that the QR iteration is equivalent to orthogonal iteration with $Q_0 = I$ (\Cref{sec:basic-qr}). The shifted QR iteration is likewise equivalent to a \emph{shifted} orthogonal iteration. With $T_0 = H$, the first two steps read $H - \mu_1I = U_1R_1$, $T_1 = R_1U_1 + \mu_1I$, and $T_1 - \mu_2I = U_2R_2$. Hence
\begin{equation}
  \begin{aligned}
    U_1U_2R_2R_1 &= U_1(T_1 - \mu_2I)R_1 = U_1(R_1U_1 + \mu_1I - \mu_2I)R_1 \\
    &= (H - \mu_1I)\bigl((H - \mu_1I) + \mu_1I - \mu_2I\bigr) \\
    &= (H - \mu_1I)(H - \mu_2I),
  \end{aligned}
\end{equation}
where the second line uses $U_1R_1 = H - \mu_1I$ twice: $U_1R_1U_1R_1 = (H - \mu_1I)^2$ and $U_1(\mu_1 - \mu_2)R_1 = (\mu_1 - \mu_2)(H - \mu_1I)$.

In the same way we can write, with $\bar{Q}_k \defeq U_1\cdots U_k$ and $\bar{R}_k \defeq R_k\cdots R_1$,
\begin{equation}
  \label{eq:shifted-qr-poly}
  (H - \mu_kI)\cdots(H - \mu_1I) = \bar{Q}_k\bar{R}_k.
\end{equation}
For the induction step, note that $T_{k-1} = \bar{Q}_{k-1}^\ast H\bar{Q}_{k-1}$ by the similarity relation of \Cref{sec:shifted-qr}, so
\begin{equation}
  \bar{Q}_k\bar{R}_k = \bar{Q}_{k-1}(T_{k-1} - \mu_kI)\bar{R}_{k-1} = (H - \mu_kI)\bar{Q}_{k-1}\bar{R}_{k-1}.
\end{equation}
$\bar{Q}_k$ is unitary and $\bar{R}_k$ is upper triangular as a product of upper triangular matrices, so \eqref{eq:shifted-qr-poly} is a QR factorization of the matrix polynomial $p_k(H) = (H - \mu_kI)\cdots(H - \mu_1I)$. The factors commute, so their order does not matter.

Compared with the unshifted case, where $\bar{Q}_k$ is the Q factor of $H^k$, the power $H^k$ is replaced by $p_k(H)$, which multiplies the eigenvector component for $\lambda_j$ by $p_k(\lambda_j) = \prod_{i=1}^{k}(\lambda_j - \mu_i)$ instead of $\lambda_j^k$. If the shifts are close to an eigenvalue, $p_k$ is tiny there, which is what accelerates convergence. Moreover, $T_k = \bar{Q}_k^\ast H\bar{Q}_k$ shows how to track the eigenvectors: $\bar{Q}_k$ accumulates all the transformations, and once $T_k$ has converged to upper triangular form, $Q\bar{Q}_k$ contains the Schur vectors of $A$.

\paragraph{What about inverse iteration?} Inverting \eqref{eq:shifted-qr-poly} gives $p_k(H)^{-1} = \bar{R}_k^{-1}\bar{Q}_k^\ast$, and taking the adjoint,
\begin{equation}
  \label{eq:shifted-qr-inverse}
  \bar{Q}_k\bar{R}_k^{-\ast} = (H - \mu_kI)^{-\ast}\cdots(H - \mu_1I)^{-\ast},
\end{equation}
where $X^{-\ast} \defeq (X^{-1})^\ast$. Inversion and the adjoint both reverse the order of a product, but the factors on the right commute, so the order does not matter. The right-hand side is exactly the operator of shifted inverse iteration (\Cref{alg:inverse}), applied to $H^\ast$ with shifts $\bar{\mu}_i$ ($H^\ast$ has the eigenvalues $\bar{\lambda}_j$).

The difference to \eqref{eq:shifted-qr-poly} is the triangular factor: $\bar{R}_k^{-1}$ is upper triangular, so $\bar{R}_k^{-\ast}$ is \emph{lower} triangular. For an upper triangular factor, the first columns of the product only involve the first columns of $\bar{Q}_k$; for a lower triangular factor, the last columns only involve the last columns of $\bar{Q}_k$. In particular, the last column of $\bar{R}_k^{-\ast}$ is $(\bar{R}_k^{-\ast})_{mm}\vec{e}_m$ with $(\bar{R}_k^{-\ast})_{mm} = 1/\overline{(\bar{R}_k)_{mm}}$, so multiplying \eqref{eq:shifted-qr-inverse} by $\vec{e}_m$ gives
\begin{equation}
  \bar{\vec{q}}_m = \overline{(\bar{R}_k)_{mm}}\;(H - \mu_kI)^{-\ast}\cdots(H - \mu_1I)^{-\ast}\vec{e}_m,
\end{equation}
where $\bar{\vec{q}}_m$ is the last column of $\bar{Q}_k$. So the same matrix $\bar{Q}_k$ carries out two iterations at once:
\begin{center}
  \begin{tabular}{lll}
    \toprule
    factorization & triangular factor & columns of $\bar{Q}_k$ \\
    \midrule
    $p_k(H) = \bar{Q}_k\bar{R}_k$ & upper & first columns: orthogonal (power) iteration \\
    $p_k(H)^{-\ast} = \bar{Q}_k\bar{R}_k^{-\ast}$ & lower & last columns: shifted inverse iteration \\
    \bottomrule
  \end{tabular}
\end{center}
The last column is shifted inverse iteration started from $\vec{e}_m$; for a single step this is the argument of~\cite[Section 4.4.4]{Demmel}. In short, \emph{the QR iteration is both power and inverse power iteration}. This explains why the bottom right corner of $T_k$ converges first, and why the shift should target $\lambda_m$ (\Cref{sec:choosing-shift}): a shift close to an eigenvalue makes inverse iteration converge fast (\Cref{thm:inverse-iter}).

\section{(SUPPLEMENT) Further Remarks}

\paragraph{The Rayleigh quotient shift.} The board leaves the concrete choice of $\mu_k \approx \lambda_m$ open. The natural estimate is the bottom right entry, $\mu_k = t_{mm}^{(k-1)} = \vec{e}_m^\ast T_{k-1}\vec{e}_m$, which is the Rayleigh quotient of $T_{k-1}$ at $\vec{e}_m$. This connects to Rayleigh quotient iteration (\Cref{alg:rqi}): by the inverse iteration interpretation above, the last column of $\bar{Q}_k$ is inverse iteration with the shift taken from the current Rayleigh quotient. Demmel~\cite[Section 4.4.4]{Demmel} explains the resulting local \emph{quadratic} convergence as follows. If $\eta = \norm{T_{k-1}(m, 1{:}m-1)}/\norm{A} \ll 1$, then setting the last row's off-diagonal part to zero perturbs an eigenvalue by $O(\eta\norm{A})$ (if it is well separated), so $\abs{\lambda_m - t_{mm}} = O(\eta\norm{A})$. The next step shrinks the last row by the factor $\abs{\lambda_m - \mu}/\min_{j \ne m}\abs{\lambda_j - \mu} = O(\eta)$, giving $O(\eta^2)$.

\paragraph{The Wilkinson shift.} The Rayleigh quotient shift can stall. For $T_0 = \left[\begin{smallmatrix} 0 & 1 \\ 1 & 0 \end{smallmatrix}\right]$ we get $\mu_1 = t_{22} = 0$; $T_0$ is itself orthogonal, so its QR factorization is $U_1 = T_0$, $R_1 = I$, and $T_1 = R_1U_1 = T_0$. Nothing changes, because the eigenvalues $\pm 1$ are equally far from the shift. The \emph{Wilkinson shift} instead takes $\mu_k$ to be the eigenvalue of the trailing $2 \times 2$ block of $T_{k-1}$ that is closer to $t_{mm}$; in the example this is $\pm 1$, an exact eigenvalue, and one step suffices. For symmetric tridiagonal matrices, QR iteration with the Wilkinson shift is globally and at least linearly convergent, and asymptotically cubically convergent for almost all matrices~\cite[Theorem 5.8]{Demmel}.

\paragraph{When is a subdiagonal entry $\approx 0$?} Since $T_k$ is obtained from $A$ by unitary similarity transformations in floating point, it already carries rounding errors of size $O(\epsm)\norm{A}$. A subdiagonal entry below this size may as well be zero, and setting it to zero is a backward stable perturbation~\cite[Section 4.4.5]{Demmel}. The iteration ends when all diagonal blocks are $1 \times 1$ (or $2 \times 2$ for complex conjugate pairs of a real matrix).

\paragraph{Why the implicit Q theorem matters for shifts.} With a shift $\mu$, one QR step factors $T_{k-1} - \mu I = U_kR_k$ and sets $T_k = R_kU_k + \mu I = U_k^\ast T_{k-1}U_k$. By \Cref{thm:implicit-q}, $U_k$ is determined by its first column, which is parallel to the first column of $T_{k-1} - \mu I$, i.e., to $(t_{11} - \mu, t_{21}, 0, \dots, 0)^\T$. So $G_1$ is computed from these two numbers, and the rest of the step is the bulge chasing of \Cref{alg:bulge-chasing} on $T_{k-1}$ itself; $T_{k-1} - \mu I$ is never formed. This is the \emph{implicitly shifted} QR algorithm~\cite[Section 4.4.8]{Demmel}. For real matrices with complex conjugate eigenvalues, two such steps with shifts $\mu$ and $\bar{\mu}$ can be combined into a \emph{double-shift} step that stays in real arithmetic (the Francis step)~\cite[Section 4.4.8]{Demmel}.


\paragraph{The Krylov view.} Why should the first column determine everything? From $AQ = QH$, column $i$ gives $h_{i+1,i}\vec{q}_{i+1} = A\vec{q}_i - \sum_{j \le i}h_{ji}\vec{q}_j$. By induction, using $h_{i+1,i} \ne 0$,
\begin{equation}
  \spn(\vec{q}_1, \dots, \vec{q}_k) = \spn(\vec{q}_1, A\vec{q}_1, \dots, A^{k-1}\vec{q}_1) \qquad \text{for } k = 1, \dots, m.
\end{equation}
The right-hand side is the Krylov subspace generated by $A$ and $\vec{q}_1$, and depends on nothing else. Equivalently, the Krylov matrix $K = [\vec{q}_1, A\vec{q}_1, \dots, A^{m-1}\vec{q}_1]$ satisfies $K = QR$ with $R = Q^\T K$ upper triangular and invertible. So $Q$ is the orthogonal factor of the QR factorization of $K$, and the implicit Q theorem follows from the uniqueness of QR up to signs~\cite[Lecture 33]{TB}. The proof above is the same argument written out column by column.

\paragraph{A direct check for the first bulge.} For the first rotation after $G_1$ one can also see directly that the two algorithms agree. Write $G_1 = \left[\begin{smallmatrix} c & -s \\ s & c \end{smallmatrix}\right]$ in rows and columns $1, 2$, and let $X = G_1^\T T_{k-1}$. The first column of $X$ is zero in rows $2, \dots, m$, and the two-sweep algorithm chooses $G_2$ from $(x_{22}, x_{32})$. Right multiplication by $G_1$ replaces column $1$ by $c\,(\text{column } 1) + s\,(\text{column } 2)$, so in $XG_1$ the bulge column has entries $s\,(x_{22}, x_{32})$ in rows $2, 3$. This vector is parallel to $(x_{22}, x_{32})$, so bulge chasing chooses the same $G_2$ up to sign (if $s = 0$, then $(T_{k-1})_{21} = 0$ and $T_{k-1}$ was reduced). Continuing this argument down the matrix is possible but tedious; the implicit Q theorem handles all rotations at once.
